\begin{theorem}[conductor-supported ambient differential]
\hypertarget{tag-03IK}{}\label{wbII:prop:conductor}
\StatementRecord{wbII:prop:conductor}{proved}{03IK}{\SeriesRef{wbII:eq:W:5.3}, \SeriesRef{wbII:eq:W:5.5}, \SeriesRef{wbII:eq:W:5.4}}
\label{wbII:prop:conductor-differential}
\ In the setting \eqref{wbII:eq:W:5.3}, let $A$ be a line bundle on $S$ with
$\eta^*A\simeq\mathcal O_{\widetilde S}$, and choose a nowhere-vanishing
section $e$ of $\eta^*A$.  There is a unique nonzero
\[
s\in H^0(S,\Omega^1_{\mathcal X}|_S\otimes A)
\tag{W.6}\hypertarget{tag-03HZ}{}\label{wbII:eq:W:5.6}
\]
such that
\[
\eta^*s=\eta^*(dg|_S)\otimes e,\qquad s|_D=0.
\tag{W.7}\hypertarget{tag-03I0}{}\label{wbII:eq:W:5.7}
\]
Its image $\rho(s)$ in $\Omega^1_S\otimes A$ is supported on $D$ and
therefore torsion; its image in
$\widetilde\Omega^1_S\otimes A$ is zero.  Moreover:
\begin{enumerate}[label=(\alph*)]
\item $\rho(s)=0$ if and only if $e=\eta^*\sigma$ for a nowhere-vanishing
$\sigma\in H^0(S,A)$;
\item if $S$ is compact and connected and $A\not\simeq\mathcal O_S$, then
$H^0(S,A)=0$ and $s$ does not come from the conormal term in \CompendiumRef{II}{wbII:eq:W:5.1}.
\end{enumerate}
\end{theorem}



\begin{proof}

\begin{enumerate}[label={$\langle1\rangle$\arabic*.},leftmargin=*]
\item In the double model, $d(xy)=y\,dx+x\,dy$ has coefficients in the reduced
ideal $(x,y)$ of $D$.  In the triple model,
\[
d(xyz)=yz\,dx+xz\,dy+xy\,dz
\]
has coefficients in $(xy,xz,yz)$, the reduced ideal of the union of the
three double curves.  Unit factors in \eqref{wbII:eq:W:5.3} do not change this after
restriction to $S$.  Hence $\eta^*(dg|_S)$ restricts to zero on
$\widetilde D$.  Tensor \eqref{wbII:eq:W:5.5} by the locally free sheaf
$E=\Omega^1_{\mathcal X}|_S\otimes A$.  The compatible pair
\[
\bigl(\eta^*(dg|_S)\otimes e,0\bigr)
\]
therefore descends uniquely to $s$.  On the smooth locus of $S$, $dg|_S$
is a nonzero conormal vector and $e$ is nowhere zero, so $s\ne0$.

\item On $S\setminus D$, the normalization is an isomorphism and
$s=dg|_S\otimes e$.  By \eqref{wbII:eq:W:5.4} its intrinsic image vanishes there.  Thus
$\rho(s)$ is supported on $D$ and lies in the torsion subsheaf; its image
modulo torsion is zero.

\item If $\rho(s)=0$, exactness of \eqref{wbII:eq:W:5.4} after tensoring by $A$ gives
$s=dg|_S\otimes\sigma$ for a section $\sigma$ of $A$.  Pulling back and
comparing with \eqref{wbII:eq:W:5.7} on the dense open set where $dg\ne0$ gives
$\eta^*\sigma=e$ there, hence everywhere by holomorphic continuation.
Since $e$ never vanishes and $\eta$ is surjective, $\sigma$ never vanishes
and trivializes $A$.  The converse is immediate.

\item Finally suppose $S$ is compact and connected and $\sigma\in H^0(S,A)$.  On each compact
connected component of $\widetilde S$, write
$\eta^*\sigma=c_i e$; the holomorphic function $c_i$ is constant.  If one
$c_i$ is zero and another is nonzero, connectedness of the finite union of
irreducible components of $S$ provides two such components meeting at a
point.  The value of $\sigma$ at that point would be both zero and nonzero,
a contradiction.  Hence either every $c_i=0$, so $\sigma=0$, or every
$c_i\ne0$, so $\sigma$ is nowhere zero and $A$ is trivial.  When $A$ is
nontrivial only the first case occurs.  Since
$N^*_{S/\mathcal X}\simeq\mathcal O_S$, the conormal global-section space is
$H^0(S,A)=0$, while $s\ne0$.

\end{enumerate}
\end{proof}